\begin{textsample}{POS Dim 1 – ce – Score 72.00 – ce/ce\_inf\_e\_ef\_11.txt}  \label{ex:f1_pos_042}
2.3.
Concepções Consideramos importante \textbf{que} o currículo a ser implementado mantenha coerência com os propósitos do DCRC e efetivamente contribua para a construção da escola necessária para o \textbf{século} XXI.
Neste sentido, entendemos como fundamental apresentar as principais concepções pedagógicas adotadas neste documento : Seres Humanos — são seres históricos e sociais \textbf{que} pensam, raciocinam, deduzem e abstraem, além de pessoas \textbf{que} sentem, se emocionam, desejam, imaginam e se sensibilizam.
Pessoas \textbf{que}, em suas relações sociais com o outro e com o mundo, formam -se integrais e integrados(as), participativos(as), ousados(as), reflexivos(as), críticos(as), autônomos(as), livres de preconceitos, criativos(as), curiosos(as), investigadores(as), solidários(as), cooperativos(as) e construtores de sua realidade.
Sociedade — é complexa e está em permanente processo de transformação.
O currículo a ser desenvolvido deve contribuir para a construção de \textbf{uma} sociedade democrática, socialmente \textbf{justa}, inclusiva, humana e solidária.
Com base em Castells ( 1999 ), o currículo a ser desenvolvido integra -se à concepção de sociedade da informação, contextualizando -se nesse período histórico, caracterizado pela revolução tecnológica \textbf{que} se consubstancia pelas tecnologias digitais de informação e de comunicação.
Dessa forma, pensamos em \textbf{fomentar} a permanente reflexão crítica, para \textbf{que} seja favorecida a utilização de critérios socioculturais e éticos, valorizar interações \textbf{face} a \textbf{face} e, de forma coerente, fortalecer \textbf{métodos} e estratégias para a construção de \textbf{uma} sociedade da informação global e \textbf{justa}.
Educação — constitui \textbf{um} processo consciente de livre adesão dos \textbf{sujeitos}, cuja ação educativa da escola ( \textbf{enquanto} processo formal ) cumpre \textbf{uma} função social não só de transmitir saberes \textbf{historicamente} acumulados, mas, essencialmente, de construir/reconstruir o conhecimento na perspectiva da formação de indivíduos éticos, responsáveis, comprometidos social e politicamente, integrados no tempo, no espaço, na sociedade.
Currículo — expressão do seu tempo, o currículo é “ \textbf{um} conjunto de decisões sobre o projeto formativo de \textbf{homem}, envolvendo valores sociais e culturais, interesses e aspirações pessoais e coletivos ” ( CEARÁ, p. 9 ).
É também compreendido como : > “ Projeto educacional planejado e desenvolvido a partir de \textbf{uma} seleção da cultura e das experiências das quais deseja -se \textbf{que} as novas gerações participem, a fim de socializá -las e capacitá -las para ser cidadãos e cidadãs solidários, responsáveis e democráticos ( SANTOMÉ, 1998, 95 ). ” Essa concepção rompe, portanto, com o sentido tradicionalmente compreendido de “ rol de \textbf{disciplinas} e conteúdos ”.
Conhecimento — é o resultado de \textbf{um} processo interativo, na proporção em \textbf{que} o \textbf{sujeito} se relaciona com o objeto, modificando -o e sendo por ele cognitivamente modificado.
Conforme Freire ( 2009 ), “ o conhecimento emerge através da invenção e reinvenção, através de \textbf{um} questionamento inquieto, impaciente, continuado e esperançoso de \textbf{homens} no mundo, com o mundo e entre si. ” ( \textbf{FREIRE}, 2009, p. 33 ) Alfabetização / Letramento — “ alfabetizar na perspectiva do letramento é possibilitar \textbf{que} os alunos aprendam a Língua Portuguesa, usufruindo e descobrindo os sentidos significados das práticas socioculturais de oralidade, leitura e escrita ; a apropriação do sistema de escrita alfabética e ortográfica no uso das práticas socioculturais e nos procedimentos da linguagem envolvidos : emissão, recepção e sentido.
Acreditamos \textbf{que} os alunos aprendem a ler e escrever, com melhor qualidade, quando o professor enxerga a dimensão da alfabetização e do letramento de modo indissociável ” ( CEARÁ, 2014, p. 12 ).
Numeramento — também entendido como matematizar, no \textbf{dizer} de Simonetti, não é simplesmente aprender a contar, calcular, ler e escrever números/algarismos, assim como alfabetizar não é apenas saber ler e escrever.
Matematizar é ir além, é numeralizar, como define Nunes e Bryant ( 1997 ) : > Ser numeralizado significa pensar matematicamente sobre situações.
Para pensar matematicamente, precisamos conhecer sistemas matemáticos de representação \textbf{que} utilizaremos como ferramentas.
Estes sistemas devem ter sentido, ou seja, devem estar relacionados às situações nas quais podem ser usados.
Precisamos ser capazes de entender a lógica dessas situações, as invariáveis, para \textbf{que} \textbf{possamos} escolher as formas apropriadas de matemática.
Desse modo, não é suficiente aprender procedimentos ; é necessário transformar esses procedimentos em ferramentas de pensamento ( NUNES E BRYANT 1997, p. 31 ).
Simonetti reafirma \textbf{que}, para matematizar, é fundamental considerar a especificidade do processo de letramento matemático e a especificidade do processo de alfabetização matemática.
Nesse sentido, faz \textbf{uma} analogia com a citação de Soares ( 2003c, p. 9,12 ) sobre alfabetização e letramento.
A autora informa \textbf{que} : > [... ] “ defender as questões específicas do processo de alfabetização não significa separar alfabetização do processo de letramento ” [... ] idem para a “ alfabetização matemática ” e o “ letramento matemático ”. [... ] “ a entrada da criança no mundo da escrita se dá simultaneamente pela aquisição do sistema convencional de escrita - a alfabetização, e pelo desenvolvimento de habilidades de uso desse sistema em atividades de leitura e escrita, nas práticas sociais \textbf{que} envolvem a língua escrita – o letramento ” [... ] idem a entrada da criança no mundo dos números se dá simultaneamente pela aquisição do conceito de número e do sistema numérico decimal convencional - a “ alfabetização matemática ”, e pelo desenvolvimento e uso desses conceitos em atividades de situações matemáticas cotidianas ( situações de quantidades, cálculos, operações, medidas ) nas práticas sociais \textbf{que} envolvem a matemática – o “ letramento matemático ”. [... ] “ a alfabetização se desenvolve no contexto de e por meio de práticas sociais de leitura e escrita, isto é, através de atividades de letramento, e este, por sua vez, só pode desenvolver -se no contexto da e por meio da aprendizagem das relações fonema-grafema, isto é, em dependência da alfabetização ” [... ] idem, também, para a “ alfabetização matemática ” e o “ letramento matemático ” \textbf{que} se desenvolvem no contexto de e por meio de práticas culturais de matemática, isto é, do “ letramento matemático ”, e este, por sua vez, só pode desenvolver -se no contexto em dependência da “ alfabetização matemática ” e as suas múltiplas “ facetas ” : sistema de numeração, operações aritméticas, medidas, geometria e tratamento de informações.
Matematizar é Numeralização – Aritmetização - Geometrização na perspectiva do letramento matemático : em/no uso das práticas sociais, culturais da matemática e os procedimentos matemáticos nelas envolvidos ( SOARES, 2003c, p. 9,12 ).
Simonetti afirma \textbf{que} a numeralização e aritmetização são a compreensão do conceito de Número, do Sistema de Numeração Decimal e das Operações Aritméticas.
Assim, como para se alfabetizar o \textbf{aprendiz} precisa compreender como o Sistema de Escrita Alfabética-SEA se organiza, para matematizar -se, ele precisa compreender como o Sistema de Numeração Decimal-SND se organiza.
Aprender sobre o \textbf{que} os números representam ( \textbf{notam} ) e aprender como os números se organizam nas notações para representar quantidades.
Desta feita, essa posição faz \textbf{uma} analogia com a citação de Artur Moraes ( 2004 ), sobre alfabetização e letramento.
Segundo ele, para compreender \textbf{que} a escrita matemática representa quantidades considerando as regras/convenções do SND, a/o \textbf{aprendiz}, necessariamente, precisará entender o conceito de número, entender a lógica do SND.
Para compreender como o SND funciona a/o \textbf{aprendiz} precisa descobrir os “ segredos ” do sistema numérico em reflexão metacognitiva : o \textbf{que} numerais/algarismos representam e como representam : 1.
O QUE OS NUMERAIS REPRESENTAM : o \textbf{aprendiz} precisa \textbf{desvendar} \textbf{que} a escrita matemática representa quantidades. 2.
COMO OS NÚMEROS SE ORGANIZAM : o \textbf{aprendiz} precisa \textbf{desvendar} o sistema de numeração decimal, como números/algarismos/algoritmos se organizam para representar as quantidades. 3.
ATIVIDADES ESTRUTURANTES COMO ATIVIDADES DIDÁTICAS — o \textbf{aprendiz} precisa de atividades \textbf{que}, em reflexão de metacognição “ meta-numérica ”, desenvolvem a compreensão do NÚMERO-SND-OPERAÇÕES ARITMÉTICAS.
Criança — “ \textbf{sujeito} histórico e de direitos, \textbf{que} interage, brinca, imagina, fantasia, deseja, aprende, observa, experimenta, narra, questiona e constrói sentidos sobre a natureza e a sociedade, produzindo cultura ” ( DCNEI — Resolução CNE/CEB no. 05/09, artigo 4o ). \textbf{Salientamos} a compreensão de \textbf{que} a criança deve ser curiosa, fazendo muitas perguntas ; ter infinitas possibilidades de criação ; explorar o mundo ; ter muito \textbf{amor}, atenção, cuidado e segurança.
Cuidar e \textbf{Educar} — o cuidar > “ é entendido de \textbf{uma} maneira restrita, referindo -se apenas à segurança física, alimentação, sono etc., \textbf{enquanto} \textbf{educar} muitas vezes é relacionado somente a conteúdos escolares. \textbf{Contudo}, as DCNEI superam essa \textbf{divisão} e nos \textbf{lembram} \textbf{que} a educação e o cuidado das crianças acontecem de maneira indissociável.
Nessa perspectiva, cuidar é bem mais do \textbf{que} atenção aos aspectos físicos e \textbf{educar} é muito mais do \textbf{que} garantir à criança acesso a conhecimentos. \textbf{Educar} e cuidar das crianças, nessa faixa de idade, é, dentre outras coisas : > > - atender às suas necessidades, oferecendo -lhes condições de se sentir confortável, em relação a sono, fome, sede, higiene, dor etc. ; > - acolher seus afetos e alimentar sua curiosidade e expressividade ; > - dar -lhes condições para explorar o ambiente e construir sentidos pessoais, sobre o mundo e sobre si, apropriando -se de formas de agir, sentir e pensar existentes em sua cultura. > > Portanto, refletir sobre essa concepção cria novas possibilidades para a Educação Infantil e requer novas atitudes por parte dos professores ” ( CEARÁ, Orientações Curriculares para a Educação Infantil, 2011 ).
Infâncias — A concepção de infância não é \textbf{uma} construção linear e existe em diferentes contextos.
Há multiplicidade de infâncias na \textbf{contemporaneidade}, mudando de acordo com o tempo e com os variados contextos sociais, econômicos, geográficos, políticos e também de acordo com as particularidades de cada indivíduo, ou seja, dentro de \textbf{uma} sociedade, as crianças \textbf{vivem} em diferentes contextos.
Desse modo, a infância é \textbf{vista} como \textbf{uma} construção social e histórica e acena sobre a importância de \textbf{que} o ambiente social, histórico e cultural \textbf{influencia} na maneira de \textbf{viver} as infâncias, ou seja, concepção \textbf{que} respeita e compreende a criança na sua \textbf{singularidade}, história e cultura ( CEARÁ, Orientações Curriculares para a Educação Infantil, 2011 ).
A brincadeira / O brincar — “ A brincadeira constitui \textbf{uma} estratégia das mais valiosas na Educação Infantil, devendo, como apontam as DCNEI, constituir a base do trabalho pedagógico ”.
Brincar é o principal modo de expressão das crianças, a ferramenta por \textbf{excelência} para elas revolucionarem seu desenvolvimento e criarem cultura.
Nas brincadeiras \textbf{que} fazem com outras crianças, com adultos, ou mesmo sozinhas, as crianças têm oportunidade para explorar o mundo, organizar seu pensamento, trabalhar seus afetos, ter iniciativa em cada situação ( CEARÁ, 2011, p. 20 ).
Adolescências — considerado como \textbf{um} fenômeno de caracterização cultural e construção social.
Está intrinsecamente vinculado à transformação e à subjetividade do desenvolvimento do ser humano, conforme cada sociedade.
Diferente da infância e da idade adulta, é \textbf{uma} fase em \textbf{que} ocorrem transformações próprias do desenvolvimento físico, psicológico, biológico e cognitivo, sendo considerado \textbf{um} período de preparação para os papéis sociais culturalmente atribuídos aos adultos.
É, sobretudo, \textbf{um} tempo em \textbf{que} se buscam respostas, em \textbf{que} se quer explorar o mundo, em meio a conflitos, inseguranças e dúvidas.
Para a neurociência, é a fase do desenvolvimento em \textbf{que} há \textbf{um} desequilíbrio entre a maturação da parte racional ( cérebro ) e a parte emocional. “ O \textbf{lado} emocional amadurece mais rapidamente \textbf{que} o racional e isso provoca, num período entre os 16 e os 18 anos, \textbf{um} desequilíbrio em \textbf{que} o comportamento é mais emocional do \textbf{que} racional ” ( VIELLE, 2018 ).
Deste modo, frente a todas essas constatações, é essencial compreender como os adolescentes se comportam na vivência das transformações naturais do período, ao mesmo tempo em \textbf{que} sofrem a \textbf{influência} das relações sociais, das condições de vida, dos valores sociais presentes na cultura do contexto em \textbf{que} \textbf{vivem}.
As adolescências, \textbf{que} decorrem desses fatores, requerem políticas públicas compatíveis com as diferentes realidades.
Escola — é \textbf{um} espaço de interação e deve ser \textbf{uma} instituição inovadora, democrática, inclusiva, crítica, \textbf{um} polo cultural da comunidade, aberta às mudanças e à cultura digital.
É \textbf{uma} instituição educativa \textbf{que} favorece ao desenvolvimento integral de alunas e alunos ; compromete -se com seu direito de aprender ; promove o desenvolvimento de competências, dentre as quais, aprender a aprender, aprender a ser, aprender a conviver, aprender a fazer, aprender a enfrentar desafios, aprender a pensar, resolver problemas.
A escola precisa permitir a/o estudante ser curiosa/curioso, criativa/criativo, solidária/solidário, cooperativa/cooperativo, a/o protagonista, ter capacidade argumentativa, voltada/voltado a construir sucesso, vivenciar direitos e deveres ; conhecer as formas de acesso e apropriação do conhecimento elaborado, de modo \textbf{que} possa utilizar as competências autonomamente ao longo de sua vida.
A escola é, principalmente, \textbf{uma} instituição \textbf{que} deve integrar -se ao seu tempo e sociedade, Espaço Escolar / Salas de Aula — o espaço escolar é público.
Lugar para onde crianças e jovens são encaminhados com a intenção de \textbf{que} sejam \textbf{educados}.
Para desenvolvimento desta missão, é fundamental \textbf{que} o espaço escolar sempre favoreça às aprendizagens.
Para tanto, precisamos, constantemente, \textbf{que}, em especial, estudantes e professores se sintam confortáveis e se apropriem do sentimento de pertença.
Constituindo espaço de convivências, o ambiente escolar em sua concretude torna -se pleno de signos, símbolos e \textbf{marcas} \textbf{que} \textbf{educam}, cuja produção e \textbf{partilha} é eminentemente pedagógica.
Na escola, onde a vida deve fluir com ampla perspectiva de objetivos a serem alcançados, a sala de aula se \textbf{configura} como \textbf{um} espaço socialmente instituído e \textbf{historicamente} construído.
Ela não é apenas \textbf{um} espaço físico, mas \textbf{um} lugar de interações onde experiências são compartilhadas e as pessoas aprendem a aprender, aprendem a fazer, aprendem a conviver, aprendem a ser, tornam -se humanas.
É a atividade desenvolvida na sala de aula \textbf{que} cria a sua especificidade, tornando -lhe efetivamente sala de aula.
Por isto, na escola, o conceito de sala de aula precisa instaurar -se em qualquer outro espaço – auditório, biblioteca, pátio, laboratórios, sanitários, diretoria, secretaria, refeitório, quadra de esporte, cozinha, projetos.
A sala de aula precisa admitir -se com formas diferentes de ser, para propiciar o trabalho compartilhado entre aluna/aluno e professora/professor ; proporcionar aulas \textbf{dialógicas}, participativas, interativas e favoráveis à criatividade ; além de estar aberta ao uso de tecnologias e à compreensão de \textbf{que} a \textbf{contemporaneidade} está interferindo nos conceitos de espaço e tempo com as possibilidades de conectividade, ubiquidade e complexidade sistêmica.
Territórios — é \textbf{um} termo da Geografia relacionado aos processos de construção e transformação do espaço geográfico.
Atualmente, o território é concebido, nas mais diversas análises e \textbf{abordagens}, como \textbf{um} espaço delimitado pelo uso de \textbf{fronteiras}, não necessariamente visíveis, consolidado a partir de relações de poder.
O território permite o entendimento das diferentes formas de apropriação do espaço, seu uso e ocupação, fazendo compreender as implicações da delimitação desses espaços por e a partir das relações de poder em determinada sociedade e em certo momento histórico.
Área, recursos, povo, poder, limites e \textbf{fronteiras} são elementos \textbf{que} definem \textbf{um} território por serem alvos diretos e indiretos dos \textbf{atores} e seus poderes.
O território é lugar de vida, do acontecer cotidiano, necessário à existência das pessoas em termos (geo)políticos, sociais, culturais, ecológicos, econômicos e políticos.
É \textbf{uma} \textbf{instância} social, faz parte da sociedade.
É expressão dela.
Assim, qualquer espaço definido e delimitado por e a partir de relações de poder, caracteriza -se como território, correspondendo ao espaço geográfico socializado, apropriado para os seus habitantes.
Professora e Professor — profissional mediadora/mediador da construção do conhecimento e do desenvolvimento de competências junto a aluna/ao aluno ; aquela/aquele \textbf{que} busca “ provocar, incentivar, disparar e possibilitar ao aluno a própria construção do conhecimento, a própria aprendizagem. ” ( MEIER e GARCIA, 2011, p. 14 ).
Tem a responsabilidade de planejar e desenvolver estratégias pedagógicas \textbf{que} promovam aprendizagens significativas.
Compreende a aprendizagem como processo contínuo, possível a todos ; busca ser fluente no uso de tecnologias digitais ; é comprometido com o direito de aprender da educanda e do \textbf{educando} e com seu aprimoramento como pessoa humana.
Profissional \textbf{que} reflete sobre seu próprio trabalho e procura \textbf{renovar} sua prática pedagógica em função da reflexão realizada.
Colaboradora/Colaborador na articulação escola/família e participante efetiva/efetivo na construção/reconstrução e execução do projeto político-pedagógico da escola.
Consciente de seu papel de \textbf{sujeito} da ação educativa e de \textbf{que} a ação de mediação da construção do conhecimento \textbf{exige} -lhe \textbf{que} tenha ciência desse conhecimento.
Aluna e aluno — estudante criativa/criativo, curiosa/curioso, investigativa/investigativo, crítica/crítico, solidária/solidário, ética/ético, cooperativa/cooperativo, responsável, argumentativa/argumentativo, resolutiva/resolutivo no sentido da tomada de decisões, dotada/dotado de iniciativa, cidadã/cidadão atuante e comprometida/comprometido com a cultura de paz e transformação da sua realidade.
Ocupa lugar \textbf{central} no processo de ensino e aprendizagem.
Constitui -se como sujeito com valores, crenças, atitudes e conhecimentos prévios, ativa/ativo na construção de seu próprio conhecimento.
Portanto, educado(a) para ser protagonista, sujeito da sua aprendizagem, entendendo \textbf{que} tem grande responsabilidade nesse processo ; também educada/educado para gerenciar emoções, negociar conflitos e buscar soluções para situações reais, \textbf{que} impactem de maneira direta ou indireta na sua própria comunidade.
Estimulada/estimulado a trazer para a sala de aula referências de sua realidade, de modo \textbf{que} os conteúdos em estudo se tornem significativos e criem/fortaleçam competências.
Aprendizagem — processo \textbf{que} avança da concepção puramente transmissiva e cumulativa de conteúdos, para \textbf{uma} concepção de aprendizagem significativa.
Assim compreendida, as/os alunas/alunos, com a cooperação da/do professora/professor ou outros \textbf{atores}, ou até mesmo sozinhas/sozinhos, constroem significado e atribuem sentido ao \textbf{que} aprendem. “ Aprender pode ser entendido como o processo de modificação do modo de agir, sentir e pensar, de cada pessoa, e \textbf{que} não pode ser atribuído à \textbf{mera} maturação orgânica, mas à experiência.
Nessa concepção, as possibilidades de aprendizagem não são resultado de processos espontâneos.
Elas requerem alguns elementos mediadores, em especial, a colaboração de diferentes parceiros na realização de alguma \textbf{tarefa}. ” ( CEARÁ, Orientações Curriculares para a Educação Infantil, 2011, p. 21 ) Ensino — ação interativa entre educadora/educador/educanda/educando/objeto do conhecimento no sentido de \textbf{que} sejam favorecidas condições \textbf{que} facilitem a construção do conhecimento pela/pelo aluna/aluno.
Como mediador dessa construção, a/o professora/professor desenvolve atividades \textbf{que} promovam questionamentos, reflexões, análises críticas, \textbf{problematizações}, levantamento de hipóteses \textbf{que} conduzam ao conhecimento \textbf{que} ela/ele pretende \textbf{que} seja aprendido. “ Ensinar é fazer \textbf{conhecido} o desconhecido. ” ( KENSKI, 2001, p. 103 ) Tecnologias na Educação — na pauta do DCRC, na competência 4, apresentada pela BNCC, as tecnologias \textbf{configuram} como suporte à linguagem, reconhecendo \textbf{que} nossos alunos representam \textbf{uma} geração cuja leitura tem como base \textbf{uma} tela plana e iluminada.
Portanto, seu entendimento se \textbf{restringe} à comunicação, reconhecendo esse recurso como \textbf{predominante} na sociedade \textbf{contemporânea} e \textbf{que}, dessa forma, deve ser contemplado pela escola.
Na competência 5, no entanto, há \textbf{um} avanço, \textbf{que} rompe a barreira do uso e parte para o campo da criação, apresentando a necessidade de \textbf{um} trabalho crítico e ético acerca do uso das tecnologias, entendendo esse uso numa perspectiva criativa.
Entendemos, \textbf{contudo}, \textbf{que} as tecnologias estão para além do uso pedagógico, pois elas acabam por criar novos modelos, estabelecendo \textbf{uma} lógica na qual criador e criatura se mesclam de tal forma, \textbf{que} \textbf{um} tem impacto no outro, ou seja, a tecnologia altera o meio e esse novo meio acaba por \textbf{demandar} novas formas e recursos.
Diante dessa constatação, não se pode compreender as tecnologias apenas como aparelhos e suas funções, mas como arcabouço tecnossocial capaz de revolucionar o \textbf{que} ainda temos como processo de ensino e aprendizagem.
Dessa maneira, as tecnologias figuram como suporte a metodologias inovadoras, abrindo caminho para \textbf{um} aprender mais significativo, com base na construção e \textbf{partilha} de saberes entre alunas/alunos, professoras/professores e comunidades científicas.
A compreensão e o domínio das tecnologias significam, na Educação Básica, \textbf{um} movimento capaz de gerar novos conteúdos em rede, além de constituir -se como componente essencial na formação dos profissionais \textbf{demandados} pelo mercado mundial. \textbf{Salientamos} duas considerações importantes : a ) a formação inicial e continuada das/dos professoras/professores e gestoras/gestores acerca do uso das tecnologias na educação continua, \textbf{configurando} -se como condição essencial para o trabalho na área, permitindo a esses profissionais \textbf{uma} aproximação ao tema de forma \textbf{teórica} e prática e, consequentemente, ampliando suas experiências no campo do ensino e da aprendizagem ; b ) faz -se necessário \textbf{que} os sistemas oficiais de ensino, \textbf{enquanto} agências de incentivo e fomento ao uso das tecnologias, mantenham parcerias com instituições de ensino superior e outras organizações com expertises na área.
Isso manterá a discussão e, consequentemente, a inovação em movimento, gerando \textbf{um} cenário de permanente estudo e evolução.
Desse modo, qualquer discussão sobre educação \textbf{que} passe pela definição de currículos e práticas \textbf{metodológicas} deve contemplar o papel das tecnologias, não apenas como ferramenta, mas como movimento \textbf{que} altera cenários e demanda novas formas de conceber ensino e aprendizagem, \textbf{uma} vez \textbf{que} tira a escola da posição de “ disseminadora ” de conteúdos e o aluno do papel de figurante.
Pensar tecnologias na educação é pensar numa educação cuja palavra de ordem é aprender a aprender.
É aprender gerando e partilhando conhecimentos.
É fazer junto, ganhando no coletivo.
É poder exercitar a \textbf{equidade}, respeitando diferentes caminhos e possibilidades \textbf{que} o aprender \textbf{exige} de cada ser.
Educação Científica — o fortalecimento da iniciação científica nas escolas públicas deve ser mais \textbf{que} \textbf{uma} possibilidade, deve ser a nova forma de a/o professora/professor atuar e interagir com os estudantes, em busca de respostas e de conhecimentos.
Segundo Demo ( 2010 ), educação científica é \textbf{vista} como \textbf{uma} das habilidades do \textbf{século} XXI, por ser este \textbf{século} marcado pela sociedade intensiva de conhecimento.
Ela é apreciada como referência fundamental de toda a trajetória de estudos básicos e superiores, com realce fundamental a tipos diversificados de ensino médio e técnico. \textbf{Hoje}, o desafio maior é produzir conhecimento e não mais apenas transmitir ( DEMO, 2010, p.15 ).
Neste sentido, produzir conhecimento não aponta apenas para o processo reconstrutivo, mecânico e técnico, mas, principalmente, para a possibilidade de cada estudante tornar -se produtor de saberes, na condição de \textbf{sujeito} \textbf{que} toma o destino em suas mãos.
Nesta perspectiva, consideramos \textbf{que}, embora a autonomia não possa ser plena, pode, por seu turno, ser muito ampliada, se soubermos aprender a manejar o conhecimento com autonomia.
Trata -se de trabalhar o desafio da autoria individual e coletiva, à medida \textbf{que} construímos oportunidades viáveis, as quais demonstram \textbf{que} o \textbf{sujeito} não \textbf{depende} de outros \textbf{que} as inventem.
Ele mesmo se dá a oportunidade, porque a cria.
Entendemos \textbf{que} a base para o funcionamento da pesquisa é o interesse das pessoas \textbf{que} participam do processo educativo e isso começa fazendo toda a diferença.
Ao tratar desta questão, Muller afirma \textbf{que} : > Para além de ser \textbf{uma} forma de trabalho qualificadora dos alunos, foi e é também \textbf{uma} forma qualificadora dos professores.
Envolver -se em pesquisa, em orientação ou organização de pesquisa, \textbf{implica} em seguir sua lógica interna, cuja principal característica é a de qualificação de todos \textbf{que} participam, \textbf{que} pensam e \textbf{que} agem segundo esse pensamento ( MÜLLER 2009, p. 09 ).
Nessa perspectiva, as redes de ensino precisam desenvolver ações em Educação Científica, induzindo gestoras/gestores, professoras/professores e estudantes a se sentirem motivados e convidados a desenvolver projetos de pesquisa científica e artístico-culturais.
Para tanto, a pesquisa deve ser trabalhada como o exercício/instrumento de (re)significação do ensino e do aprendizado, a partir da prática interdisciplinar \textbf{que} possibilita a integração curricular, com o envolvimento de todos os estudantes e professores de todas as áreas de conhecimento, a fim de interpretar, de responder ou solucionar \textbf{um} problema observado pelos \textbf{sujeitos} do contexto escolar.
Assim, integrar o currículo significa diminuir o isolamento, criando pontes entre as áreas e os objetos de conhecimento, possibilitando \textbf{que} estes ganhem significado prático e cotidiano em seu entendimento.
Possibilita, ainda, \textbf{que} as/os professoras/professores planejem de forma sistemática, contínua e integrada o trabalho \textbf{que} realizam na escola.
Esperamos \textbf{que} elas/eles possam trazer para as suas explanações e atividades \textbf{um} pouco do conhecimento amplo \textbf{que} conecta e faz sentido para quem precisa contextualizar aquilo \textbf{que} almeja apreender e como aprende ; bem como para \textbf{que} a escola, à proporção \textbf{que} reforça a sua identidade e o seu real papel na \textbf{atual} sociedade do conhecimento, integre o currículo escolar.
Para Müller ( 2009 ), a pesquisa é possível a todos, em todas as idades, em todos os níveis de conhecimento, porque é próprio do ser humano questionar o mundo e buscar respostas para essas perguntas.
Interações — “ Interações são ações compartilhadas do professor com as crianças e das crianças entre si.
Elas se dão em situações concretas por meio dos papéis \textbf{que} os parceiros vão assumindo nas situações ( por exemplo, o papel de quem pergunta, de quem explica, de quem pede colo, de quem provoca riso etc. ) ” ( Orientações Curriculares para a Educação Infantil/SEDUC, 2011, p.37 ).
É \textbf{imprescindível} considerar \textbf{que} o processo de aprendizagem \textbf{depende} diretamente de processos de interação entre os \textbf{sujeitos}.
A aprendizagem não se dá de forma isolada.
A convivência entre as pessoas favorece trocas \textbf{que} vão construindo o seu conhecimento.
Vygotsky ( 1994 ), ao destacar a importância das interações sociais, faz referência à mediação e à internalização como aspectos fundamentais para a aprendizagem.
O \textbf{autor} defende \textbf{que} a construção do conhecimento ocorre a partir de \textbf{um} intenso processo de interação entre as pessoas.
Fica, então, evidenciado \textbf{que} crianças, jovens e adultos adquirem novas formas de pensar e de agir apropriando -se de novos conhecimentos proporcionados em interações com outras pessoas.
Competências socioemocionais — A educação, com a implementação da BNCC, busca a superação da dicotomia intelecto/emoção, a fim de promover \textbf{uma} formação \textbf{que} tenha como propósito o desenvolvimento integral dos estudantes.
Atualmente, temos disponível \textbf{um} conjunto crescente de evidências científicas \textbf{que} têm contribuído para o aprofundamento desse tema.
Sabemos, por exemplo, \textbf{que} o desenvolvimento intencional e integrado das competências socioemocionais impactam diretamente a aprendizagem na escola, bem como diversos aspectos da vida, tais como saúde, escolaridade, renda, entre outros.
Esses estudos legitimam aquilo \textbf{que} todo professor já intui : as emoções estão profundamente ligadas à aprendizagem.
A formação integral do aluno é, portanto, cada vez mais \textbf{exigida}, fazendo -se necessário \textbf{que} reflitamos sobre \textbf{que} estudante desejamos formar.
A formação de pessoas autônomas, solidárias, capazes de fazer escolhas alinhadas aos seus projetos de futuro e interesses.
Pessoas capazes de articular e colocar em prática conhecimentos, valores, atitudes e habilidades importantes para a relação com os outros e \textbf{consigo} mesmo, para o estabelecimento de objetivos de vida e para o enfrentamento de desafios de maneira criativa e construtiva.
Isso \textbf{demanda} \textbf{uma} experiência escolar \textbf{que} permita o autoconhecimento e protagonismo de cada estudante.
Assumir tal \textbf{posicionamento} significa ter \textbf{um} currículo e \textbf{uma} escola \textbf{que} articulem o desenvolvimento cognitivo e socioemocional dos alunos intencionalmente. \textbf{Uma} escola \textbf{que} acolha as \textbf{singularidades} e as diversidades, \textbf{que} tenha como valor \textbf{central} o estabelecimento de relações entre professores e estudantes e entre os estudantes.
São \textbf{singularidades} e diversidades \textbf{baseadas} no respeito e na prática diária de metodologias ativas de ensino e de aprendizagem.
O estado do Ceará já vem desenvolvendo \textbf{uma} linha de ação com competências socioemocionais, cuja base é \textbf{uma} matriz \textbf{que} delimita cinco macrocompetências : - Abertura ao novo — descreve o grau com \textbf{que} \textbf{uma} pessoa se mostra criativa, sensível a diferentes estéticas, imaginativa e curiosa por ideias e outros modos de vida. - Autogestão — descreve o grau com \textbf{que} \textbf{uma} pessoa é centrada, organizada, diligente, \textbf{focada}, pontual e persistente. - Engajamento com os outros — descreve o grau com \textbf{que} \textbf{uma} pessoa tem interesse por se socializar, fazer amigos, ser gregário e experimentar o cotidiano com disposição e energia. - Amabilidade — descreve o grau com \textbf{que} \textbf{uma} pessoa é capaz de se colocar no lugar do outro, sentir compaixão, acreditar nas boas intenções de outra pessoa e respeitá -la. - Resiliência emocional — descreve o grau com \textbf{que} \textbf{uma} pessoa é \textbf{atravessada} por sentimentos de ansiedade, raiva e insegurança quando submetida a pressões interpessoais, estresse e frustrações.
A busca pela formação integral do aluno está, portanto, iniciada, com investimento em atitudes \textbf{que} favorecem ao relacionamento inter e intrapessoal, além da preparação para o enfrentamento de situações novas.
Avaliação da Aprendizagem — componente básico do processo de ensino e aprendizagem, cujo objetivo básico é verificar se o(a) aluno(a) aprendeu o \textbf{que} foi ensinado.
Perpassa todo o ato educativo, visando, também, orientar docentes e \textbf{discentes} no sentido de promover aprendizagem significativa.
Avaliamos o tratamento didático utilizado se ele está sendo favorável à ocorrência das aprendizagens buscadas ou se precisa ser revisto.
Ela tem caráter processual, contínuo e continuado, além de precisar ser caracteristicamente \textbf{justo}.
Deve exercer as funções diagnóstica, formativa e somativa, nas quais é importante prevalecer os aspectos qualitativos sobre os quantitativos.
É, portanto, \textbf{um} processo de investigação permanente em \textbf{que} todas as atividades de \textbf{que} o(a) aluno(a) participa são consideradas, valendo a compreensão de \textbf{que} toda avaliação de aprendizagem é essencialmente diagnóstica : faz \textbf{um} diagnóstico em \textbf{que} verifica dentre o \textbf{que} foi estudado, o \textbf{que} efetivamente foi aprendido.
Deve utilizar variadas metodologias, técnicas e instrumentos \textbf{que} possibilitem verificar o progresso do aluno no decorrer dos estudos realizados.
Avaliar a aprendizagem não é simplesmente atribuir \textbf{uma} nota ou conceito ; é refletir sobre o desempenho dos(as) estudantes frente aos objetivos traçados, sempre compreendendo \textbf{que} a consequência lógica do ensinar deve ser o aprender. \textbf{Um} ensino só é eficiente quando provoca a aprendizagem de alunos e alunas.
Neste aspecto, vale ressaltar \textbf{que}, se o(a) \textbf{aprendiz} efetivamente aprende o ensinado e desenvolve competências e habilidades como \textbf{raciocínio} lógico, criatividade, criticidade, poder argumentativo, resolutividade no sentido da tomada de decisões, iniciativa, autonomia para buscar conhecimentos, etc., estará naturalmente pronto para participar com sucesso, inclusive de avaliações externas.
É essencial termos clareza de \textbf{que}, diante do \textbf{compromisso} com a formação integral do(a) estudante, o processo avaliativo deve considerar, também, o desenvolvimento das competências socioemocionais.
Para tanto, o referido processo, na instituição educacional, deve passar por significativa transformação para incluir procedimentos adequados à avaliação do(a) aluno(a) no tocante ao desenvolvimento dessas competências incluídas.
Para esta avaliação específica, \textbf{um} caminho pode ser adaptação do processo de observações e registros, adotado na educação infantil.
Também a autoavaliação, como mecanismo de compartilhamento professora/professor aluna/aluno na execução da \textbf{tarefa} de avaliar, é recomendável, inclusive, como instrumento da formação em valores, dentre os quais, a responsabilidade e a honestidade.
Por fim, \textbf{lembramos} a importância da realidade social e cultural em \textbf{que} o(a) \textbf{aprendiz} \textbf{vive} ; e de \textbf{que} os(as) estudantes não podem ser penalizados pelas competências e habilidades socioemocionais ainda não desenvolvidas.
Cumpre \textbf{salientar}, à guisa de finalização deste item, \textbf{que} as concepções especificadas devem nortear o dia a dia de cada educador no desenvolvimento do trabalho \textbf{que} realiza, seja na \textbf{docência}, seja em qualquer outra atividade \textbf{que} realize – \textbf{direção} escolar, coordenação pedagógica, outros serviços na escola, gerência e serviços técnicos em órgãos do sistema de ensino.
Nossa consciência deve estar mobilizada pela convicção de \textbf{que} é tempo de construção de \textbf{uma} educação \textbf{que} se comprometa com \textbf{uma} sociedade mais humana, cooperativa, solidária e \textbf{comprometida} com a paz, para todos.

% matched lemmas: abordagem, amor, aprendiz, ator, atravessar, atual, autor, basear, central, comprometido, compromisso, configurar, conhecido, conseguir, contemporaneidade, contemporâneo, contudo, demandar, depender, desvendar, dialógico, direção, discente, disciplina, divisão, dizer, docência, educar, enquanto, equidade, excelência, exigir, face, focar, fomentar, freire, fronteira, historicamente, hoje, homem, implicar, imprescindível, influenciar, influência, instância, justo, lado, lembrar, marca, mero, metodológico, método, notar, partilha, posicionamento, possar, predominante, problematização, que, raciocínio, renovar, restringir, salientar, singularidade, sujeito, século, tarefa, teórico, um, ver, viver
\end{textsample}
